\documentclass[11pt]{article}
\usepackage[T1]{fontenc}
\usepackage[utf8]{inputenc}
\usepackage[margin=1in]{geometry}
\usepackage{amsmath,amssymb,amsthm,hyperref}
\hypersetup{colorlinks=true,urlcolor=blue,linkcolor=black}
\newtheorem{theorem}{Theorem}
\newtheorem{lemma}[theorem]{Lemma}
\newtheorem{corollary}[theorem]{Corollary}
\newcommand{\Q}{\mathbb Q}
\newcommand{\K}{K}
\DeclareMathOperator{\QSym}{QSym}
\title{Forbidden subposets for fair series-parallel orders}
\author{Wentao Li}
\date{November 20, 2026}
\begin{document}\maketitle

An edge-decorated finite poset has weak or strict covers. The class in \href{https://ajc.maths.uq.edu.au/pdf/90/ajc_v90_p341.pdf}{Albertin--Aval--Mlodecki, Definition 4.1} starts with a singleton and allows disjoint union and ordinal sum with every new cover of the same color. Remark 4.2 asks for a minor-free description, without defining a minor operation. Use the two containment conventions of \href{https://doi.org/10.5802/alco.273}{Aval--Djenabou--McNamara, Section 4.5}: ignore colors in an ordinary induced suborder and retain colors on convex suborders, which inherit their covers.
\begin{theorem}A nonempty decorated poset is fair series-parallel if and only if it avoids the uncolored induced $N$ and the two decorated convex forks with one weak and one strict edge, one fork pointing up and one pointing down.\end{theorem}
Here $N$ has exactly the comparisons $a<b$, $c<b$, $c<d$ between distinct elements. A mixed fork consists of a pair of differently colored covers sharing their bottom endpoint or sharing their top endpoint. Three such vertices always form a convex suborder.
\begin{proof}In an ordinal sum $U\oplus V$ the new covers form the complete bipartite graph from $\max U$ to $\min V$. They are the only outgoing covers at their lower endpoints and the only incoming covers at their upper endpoints. A uniformly colored sum therefore creates no mixed fork. Unions also preserve this property. The underlying orders are series-parallel and hence induced-$N$-free by the standard characterization recalled in Remark 4.2 of the source.
Conversely, an induced-$N$-free order on more than one vertex has a nontrivial disjoint-union or ordinal-sum decomposition, by the same characterization. The summands are convex and inherit all three avoidance conditions. In the ordinal case the colors on the bipartite interface are constant along each row and along each column by the two fork tests. Since both sides are nonempty, the entire interface is monochromatic. Induction constructs the decorated poset by the allowed operations.\end{proof}
There is also a uniform convex-containment formulation. Along with the two mixed forks forbid every decorated poset which is the union of three closed intervals $[a,b],[c,b],[c,d]$ on four endpoints inducing $N$, with all decorations allowed. Every induced $N$ in a finite order has such a convex closure: its interval union is the intersection of the upset generated by $a,c$ and the downset generated by $b,d$, since $a\nleq d$. Conversely every listed obstruction contains an induced $N$. The three tests of the theorem give the finite description. If a minor operation instead allows deletion of any Hasse edge followed by taking reachability, there is no forbidden-minor characterization under that operation: deleting one of the four covers of the all-weak ordinal sum of two two-element antichains produces the excluded $N$. Every avoidance class for a transitive minor relation must be closed under that relation.

\medskip\noindent\textit{Research draft. Prepared with assistance from OpenAI Codex. Independent review pending.}\end{document}
